% !TEX root = ../main.tex

%* 默认字体
\setmainfont{Times New Roman}
\setCJKmainfont[AutoFakeBold=2.5]{SimSun}
\setCJKmonofont{WenQuanYi Zen Hei Mono}
\setmathfont{Latin Modern Math}

%* 字体加粗
\xeCJKsetup{AutoFakeBold=true}
\newCJKfontfamily{\cuheiti}[AutoFakeBold={2.5}]{Microsoft YaHei}
\newCJKfontfamily{\cusongti}[AutoFakeBold={2.5}]{SimSun}
% \newCJKfontfamily{\cukaishu}[AutoFakeBold={2.5}]{SimKai}
\newCJKfontfamily{\cufangsong}[AutoFakeBold={2.5}]{FZFangSong-Z02}
\newCJKfontfamily{\cuyahei}[AutoFakeBold={2.5}]{Microsoft YaHei}

%* 定理样式设置
\theoremstyle{definition}
\newtheorem{definition}{定义}[section]
\newtheorem{example}[definition]{例}
\newtheorem{problem}[definition]{问题}
\newtheorem{assumption}[definition]{假设}
\newtheorem{solution}[definition]{解}
\newtheorem{theorem}[definition]{定理}
\newtheorem{lemma}[definition]{引理}
\newtheorem{corollary}[definition]{推论}
\newtheorem{conjecture}[definition]{猜想}
\newtheorem{axiom}[definition]{公理}
\newtheorem{principle}[definition]{定律}

% * 自定义字体
\DeclareCaptionFont{heismall}{\zihao{-5}\heiti}

%* 代码环境
\setmonofont{Consolas}    % 设置等宽字体，比如 Consolas
%* 默认

\lstset{
    %* 行号
    numbers             =   left,               % 行号的位置在左边
    numberstyle         =   \zihao{-5}\ttfamily\color{gray},
    showspaces          =   false,              % 是否显示空格，显示了有点乱，所以不现实了
    showstringspaces    =   false,
    %* 边框
    captionpos          =   t,                  % 这段代码的名字所呈现的位置，t指的是top上面
    frame               =   lrtb,               % 显示边框
    rulecolor           =   \color{gray!45},    % 边框颜色
    %* 字体
    basicstyle          =   \zihao{-5}\ttfamily,
    keywordstyle        =   \color{blue},
    keywordstyle        =   [2] \color{red},
    stringstyle         =   \color{magenta},
    commentstyle        =   \color{teal}\ttfamily,
    %* 排版
    breaklines          =   true,               % 自动换行，建议不要写太长的行
    columns             =   fixed,              % 如果不加这一句，字间距就不固定，很丑，必须加
    basewidth           =   0.5em,
    tabsize             =   4,
    %* 背景
    backgroundcolor     =   \color{gray!5},     % 代码背景色
    %* 其他
    extendedchars       =   true,               % 允许使用扩展字符（如中文）
    inputencoding       =   utf8,               % 代码文件编码
    upquote             =   true,               % 修复反引号显示问题
    xleftmargin         =   1em,                % 代码块左边距
    xrightmargin        =   1em,                % 代码块右边距
    aboveskip           =   1em,                % 代码块上方间距
    belowskip           =   1em                 % 代码块下方间距
}

%* 表格标题格式
% \captionsetup[table]{
%     font={small,  bf},        % 设置字号（small 约等于小五）和加粗
%     labelsep=quad,            % 分隔符：编号和文字间留出一个空格
%     labelfont=bf,             % 编号（如“表 1”）加粗
%     textfont={bf},            % 标题文字本身不加粗
%     justification=centering   % 居中显示
% }
% \captionsetup[table]{
%     font=heismall,            % 全体（标签+文字）使用小五黑体
%     justification=centering   % 强制居中
% }

%* 四级标题
% 设置编号深度到第4层 (paragraph)
\setcounter{secnumdepth}{4}
% 设置目录深度到第4层
\setcounter{tocdepth}{4}

\titleformat{\paragraph}[hang]{\normalfont\zihao{-4}\bfseries}{\arabic{subsection}.\arabic{subsubsection}.\arabic{paragraph}}{0.5em}{}
\titlespacing*{\paragraph}{0pt}{3.25ex plus 1ex minus .2ex}{1.5ex plus .2ex}

%* 目录格式
% 调整 Section 编号与标题的间距 (默认通常是 1.5em)
\setlength{\cftsecnumwidth}{0em}
% 调整 Subsection 编号与标题的间距 (默认通常是 2.3em)
\setlength{\cftsubsecnumwidth}{1em}
% 如果有 Subsubsection，也可以一并调整
\setlength{\cftsubsubsecnumwidth}{2em}
\setlength{\cftparanumwidth}{2em}
\definecolor{mygreen}{RGB}{0, 150, 0} % 调低第二个数值可变深